% ============================================================
% Sección: Relación con la Unidad 1
% Responsable: Sánchez Vera Darwin Joel
% ============================================================
\section*{Relación con la Unidad 1}
\textbf{Transparencia de distribución.} La transparencia ante fallos busca que el usuario no note que una parte del sistema dejó de funcionar~\cite{vansteen2016}. En una CDN, si un servidor falla, otro debería atender la petición sin que el visitante lo perciba. El 8 de junio ocurrió lo contrario: el 85\,\% de la red devolvió errores~\cite{fastly2021} y los visitantes de sitios como Reddit recibieron errores HTTP 503 \textit{Service Unavailable}~\cite{medina2021thousandeyes}. Un proveedor intermedio que el usuario normalmente no ve quedó expuesto en su pantalla.

\textbf{Teorema CAP.} El teorema relaciona tres propiedades: consistencia (C, todos los nodos muestran los mismos datos), disponibilidad (A, toda petición recibe una respuesta válida) y tolerancia a particiones (P, el sistema sigue operando aunque se corte la comunicación entre nodos). Brewer aclara que elegir entre C y A solo es necesario cuando existe una partición~\cite{brewer2012}. En el incidente se perdió la disponibilidad: los servidores respondían, pero con errores~\cite{fastly2021,medina2021thousandeyes}. Sin embargo, no fue un sacrificio de diseño ante una partición: las fuentes no describen una división de la red, sino un defecto de software activado por una configuración válida~\cite{fastly2021}. Por eso el caso no encaja en el dilema clásico C frente a A; muestra que la disponibilidad también se pierde sin partición cuando todos los nodos comparten el mismo defecto.

\textbf{Tolerancia a fallos.} El mecanismo que no funcionó fue la redundancia por replicación: muchos servidores distribuidos que se respaldan entre sí para que el fallo de uno no afecte el servicio. Esta redundancia supone que los nodos fallan de forma independiente; Avižienis et al.\ explican que las copias idénticas no toleran defectos de desarrollo persistentes, que requieren diversidad de diseño~\cite{avizienis2004}. En Fastly, todos los nodos ejecutaban el software desplegado desde el 12 de mayo con el mismo defecto latente, y la configuración de un cliente lo activó a la vez en el 85\,\% de la red~\cite{fastly2021}: un fallo de modo común. La detección sí funcionó (09:48 UTC), pero identificar la configuración tomó 40 minutos (10:27 UTC)~\cite{fastly2021}. Esto se relaciona con dos causas: las pruebas no detectaron el defecto y la configuración de un solo cliente no quedó aislada del resto; por ello Fastly se propuso aprovechar el aislamiento de WebAssembly y Compute@Edge~\cite{fastly2021}.
